%% kapa: exam=1 round=20 year=2009 period=2 form=A subject=appraisal_law subject_ko=부동산관계법규 seq=10 number=50
%% source: Q-Net 감정평가사 기출문제 (cst00309) artlSeq=5024178 file=extracted/1008/5024178/2009년 제20회 감정평가사 1차시험 기출문제/2교시.hwp
%% conversion: kordoc 4.24.51 -> tools/gen1.py (auto); review=proofread-sonnet; flags=none
\begin{question}{50}
국토의 계획 및 이용에 관한 법령상 도시관리계획 결정의 효력에 관한 설명으로 옳지 \uline{않은} 것은?
\begin{choices}
  \item 도시관리계획 결정은 고시가 된 날부터 5일 후에 그 효력이 발생한다.
  \item 도시관리계획 결정 당시 허가를 받아 이미 사업이나 공사에 착수한 자도 그 사업이나 공사를 계속하기 위해서는 새로운 허가를 받아야 한다.
  \item 도시관리계획 결정 효력의 구체적ㆍ개별적인 범위는 지적고시도면(지형도면)에 의하여 확정된다는 것이 판례의 입장이다.
  \item 도시관리계획 결정의 고시일부터 2년이 되는 날까지 법에 따라 필요한 지형도면의 고시가 없는 경우에는 그 2년이 되는 날의 다음날에 그 도시관리계획 결정은 효력을 잃는다.
  \item 도시관리계획 결정이 효력을 잃게 된 때에는 대통령령으로 정하는 바에 따라 지체 없이 그 사실을 고시하여야 한다.
\end{choices}
\end{question}
\answer{2}
